\chapter{Alterações Estruturais, Restrições e Integridade Referencial}
\label{cap:alteracoes_restricoes_integridade}

Depois que a estrutura inicial do banco de dados é criada, raramente ela permanece inalterada durante todo o ciclo de vida do sistema. Novos requisitos aparecem, regras de negócio mudam, colunas precisam ser incluídas, tipos devem ser ajustados e vínculos entre tabelas precisam ser refinados. Por esse motivo, a DDL não se limita à criação inicial de objetos: ela também oferece mecanismos para evolução estrutural do banco.

Esta aula aprofunda justamente esse aspecto da implementação. O foco recai sobre alterações em tabelas já existentes, definição de restrições adicionais e decisões ligadas à integridade referencial, especialmente no tratamento de remoções e atualizações em registros relacionados.

\section{Evolução estrutural do banco}

Em um projeto acadêmico, comercial ou institucional, a base de dados é construída para representar uma realidade volátil. Como o mundo real muda constantemente, a estrutura do banco deve acompanhar essas mutações por meio da inclusão de novas colunas em tabelas legadas, da redefinição emergencial de tipos de dados, do endurecimento de regras através de novas restrições lógicas, da reconfiguração minuciosa de chaves estrangeiras e, invariavelmente, da exclusão cirúrgica de objetos que se tornaram obsoletos com o tempo.

Essa dinâmica exige domínio do comando \texttt{ALTER TABLE} e das demais operações de manutenção estrutural.

\section{ALTER TABLE}

O comando \texttt{ALTER TABLE} é um dos mais importantes da DDL. Ele permite modificar a definição de uma tabela sem recriá-la do zero. Isso é especialmente valioso em sistemas que já possuem dados armazenados.

\subsection{Adicionando colunas}

\begin{sqlcode}
ALTER TABLE alunos
ADD COLUMN email VARCHAR(120);
\end{sqlcode}

Esse comando acrescenta uma nova coluna à tabela \texttt{alunos}. Em um cenário real, isso pode ocorrer quando o sistema passa a exigir o armazenamento de contato eletrônico dos estudantes.

\subsection{Adicionando mais de uma coluna}

\begin{sqlcode}
ALTER TABLE alunos
ADD COLUMN telefone VARCHAR(20),
ADD COLUMN data_nascimento DATE;
\end{sqlcode}

Em alguns SGBDs, múltiplas alterações podem ser realizadas em um único comando. Isso é útil para migrações estruturais maiores.

\subsection{Alterando tipo de coluna}

\begin{sqlcode}
ALTER TABLE alunos
ALTER COLUMN email TYPE VARCHAR(150);
\end{sqlcode}

Esse ajuste é necessário, por exemplo, quando se percebe que o tamanho inicialmente definido foi insuficiente.

\subsection{Renomeando coluna}

\begin{sqlcode}
ALTER TABLE alunos
RENAME COLUMN curso TO curso_nome;
\end{sqlcode}

Esse tipo de alteração é mais semântica do que estrutural, mas pode melhorar a clareza do banco.

\subsection{Removendo colunas}

\begin{sqlcode}
ALTER TABLE alunos
DROP COLUMN telefone;
\end{sqlcode}

Remover colunas exige cautela, pois pode causar perda de dados e impacto em aplicações que dependem da estrutura antiga.

\begin{danger}
Antes de remover uma coluna, é essencial verificar se ela não está sendo usada por consultas, relatórios, integrações ou rotinas da aplicação.
\end{danger}

\section{DEFAULT}

A cláusula \texttt{DEFAULT} define um valor que será assumido automaticamente quando nenhum outro for informado.

\begin{sqlcode}
ALTER TABLE alunos
ADD COLUMN status VARCHAR(20) DEFAULT 'ATIVO';
\end{sqlcode}

Nesse exemplo, todo novo aluno passa a receber o status \texttt{ATIVO}, salvo se outro valor for explicitamente informado.

\subsection{DEFAULT com datas}

\begin{sqlcode}
ALTER TABLE notas
ALTER COLUMN data_avaliacao SET DEFAULT CURRENT_DATE;
\end{sqlcode}

Essa configuração é útil quando o sistema registra automaticamente a data da operação.

\subsection{Importância dos valores padrão}

A configuração diligente de valores padrão é uma técnica poderosa de engenharia de dados. Ao delegar essa responsabilidade ao SGBD, reduz-se drasticamente o número de campos esquecidos nos formulários (\textit{nulls} não intencionais), eleva-se a consistência geral da base, diminui-se o acoplamento perigoso com o código da aplicação cliente, e as regras estritas do domínio de negócio passam a viver enraizadas na própria estrutura da tabela.

\section{CHECK}

A restrição \texttt{CHECK} permite impor condições lógicas sobre valores aceitos por uma coluna ou conjunto de colunas.

\subsection{Exemplo com idade}

\begin{sqlcode}
ALTER TABLE alunos
ADD CONSTRAINT chk_idade
CHECK (idade >= 0);
\end{sqlcode}

\subsection{Exemplo com nota}

\begin{sqlcode}
ALTER TABLE notas
ADD CONSTRAINT chk_nota
CHECK (nota >= 0 AND nota <= 10);
\end{sqlcode}

\subsection{Exemplo com status}

\begin{sqlcode}
ALTER TABLE alunos
ADD CONSTRAINT chk_status
CHECK (status IN ('ATIVO', 'INATIVO', 'TRANCADO'));
\end{sqlcode}

O uso de \texttt{CHECK} é importante porque transfere parte da validação do nível da aplicação para o próprio banco de dados.

\section{UNIQUE}

A restrição \texttt{UNIQUE} impede repetição de valores em determinada coluna ou combinação de colunas.

\begin{sqlcode}
ALTER TABLE alunos
ADD CONSTRAINT uq_email UNIQUE (email);
\end{sqlcode}

Esse tipo de restrição é útil quando se deseja garantir, por exemplo, que não existam dois alunos com o mesmo e-mail institucional.

\subsection{UNIQUE composta}

\begin{sqlcode}
ALTER TABLE matriculas
ADD CONSTRAINT uq_matricula_unica
UNIQUE (aluno_id, disciplina_id, semestre);
\end{sqlcode}

Nesse caso, impede-se repetição da mesma matrícula de um aluno na mesma disciplina e no mesmo semestre.

\section{NOT NULL em colunas existentes}

Às vezes uma coluna criada inicialmente como opcional passa a ser obrigatória.

\begin{sqlcode}
ALTER TABLE alunos
ALTER COLUMN email SET NOT NULL;
\end{sqlcode}

Antes de aplicar esse tipo de alteração, é preciso garantir que não haja valores nulos já armazenados.

\section{DROP TABLE e IF EXISTS}

\begin{sqlcode}
DROP TABLE IF EXISTS temporaria;
\end{sqlcode}

Esse comando remove a tabela somente se ela existir, evitando falhas em scripts de teste e reexecução.

\subsection{Uso didático do DROP}

Em ambientes de aula e prática, é comum reiniciar o banco várias vezes. Por isso, o padrão:

\begin{sqlcode}
DROP TABLE IF EXISTS tabela;
CREATE TABLE tabela (...);
\end{sqlcode}

é bastante útil.

\section{Integridade referencial}

A integridade referencial é uma das propriedades centrais dos bancos relacionais. Ela garante que valores de chaves estrangeiras correspondam a registros válidos na tabela referenciada.

Por exemplo, se uma tabela \texttt{matricula} possui uma coluna \texttt{aluno\_id} como chave estrangeira, então não deve existir matrícula associada a um aluno inexistente.

\section{ON DELETE e ON UPDATE}

Ao definir uma chave estrangeira, é possível estabelecer o comportamento do banco quando o registro da tabela pai é alterado ou removido.

\begin{sqlcode}
CREATE TABLE matricula (
    matricula_id SERIAL PRIMARY KEY,
    aluno_id INTEGER,
    FOREIGN KEY (aluno_id)
        REFERENCES alunos(aluno_id)
        ON DELETE CASCADE
        ON UPDATE CASCADE
);
\end{sqlcode}

Essas cláusulas são extremamente relevantes porque traduzem decisões de negócio para o nível estrutural do banco.

\section{Comportamentos possíveis}

\subsection{CASCADE}

Quando o registro pai é alterado ou removido, as linhas dependentes acompanham a mudança.

\begin{sqlcode}
CREATE TABLE pedidos (
    pedido_id SERIAL PRIMARY KEY,
    cliente_id INTEGER,
    FOREIGN KEY (cliente_id)
        REFERENCES cliente(cliente_id)
        ON DELETE CASCADE
);
\end{sqlcode}

Se um cliente for removido, seus pedidos também serão excluídos.

\begin{quadro_dica}
\texttt{CASCADE} costuma ser adequado quando o registro dependente não faz sentido sem o registro principal.
\end{quadro_dica}

\subsection{SET NULL}

Quando o registro pai é removido ou alterado, a chave estrangeira da tabela filha recebe \texttt{NULL}.

\begin{sqlcode}
CREATE TABLE turma (
    turma_id SERIAL PRIMARY KEY,
    professor_id INTEGER,
    FOREIGN KEY (professor_id)
        REFERENCES professor(professor_id)
        ON DELETE SET NULL
);
\end{sqlcode}

Esse comportamento é útil quando a ausência temporária da relação não invalida o registro dependente.

\subsection{SET DEFAULT}

A chave estrangeira assume o valor padrão definido na coluna.

\begin{sqlcode}
CREATE TABLE atendimento (
    atendimento_id SERIAL PRIMARY KEY,
    funcionario_id INTEGER DEFAULT 1,
    FOREIGN KEY (funcionario_id)
        REFERENCES funcionario(funcionario_id)
        ON DELETE SET DEFAULT
);
\end{sqlcode}

\subsection{RESTRICT}

Impede a remoção ou atualização se houver registros dependentes.

\begin{sqlcode}
CREATE TABLE inscricao (
    inscricao_id SERIAL PRIMARY KEY,
    curso_id INTEGER,
    FOREIGN KEY (curso_id)
        REFERENCES curso(curso_id)
        ON DELETE RESTRICT
);
\end{sqlcode}

Nesse caso, um curso não poderá ser removido se existir inscrição vinculada.

\subsection{NO ACTION}

A cláusula \texttt{NO ACTION} é frequentemente confundida com a \texttt{RESTRICT}, pois seu objetivo principal também é abortar a operação e levantar um erro se houver registros dependentes na tabela filha. No entanto, há uma diferença arquitetural sutil e de extrema importância: a política de adiamento (\textit{deferred constraints}). Em SGBDs avançados (como o PostgreSQL), ao estipular \texttt{NO ACTION}, o desenvolvedor permite que a verificação de integridade referencial seja momentaneamente suspensa e checada apenas no exato momento do \texttt{COMMIT} final da transação, e não imediatamente a cada linha deletada.

\begin{sqlcode}
CREATE TABLE reserva (
    reserva_id SERIAL PRIMARY KEY,
    quarto_id INTEGER,
    FOREIGN KEY (quarto_id)
        REFERENCES quarto(quarto_id)
        ON DELETE NO ACTION DEFERRABLE INITIALLY DEFERRED
);
\end{sqlcode}

Esse comportamento transacional é vital para algoritmos que precisem apagar registros de chaves cruzadas (cíclicas) temporariamente dentro de um grande bloco lógico, contanto que ao final da transação todas as chaves voltem a fazer sentido.

\section{Comparando as opções}

\begin{center}
\begin{tabular}{|p{3cm}|p{8cm}|}
\hline
\textbf{Opção} & \textbf{Efeito} \\
\hline
CASCADE & Propaga remoção ou atualização aos dependentes \\
\hline
SET NULL & Define a chave estrangeira como nula \\
\hline
SET DEFAULT & Define a chave estrangeira com valor padrão \\
\hline
RESTRICT & Impede a operação caso haja dependentes \\
\hline
NO ACTION & Semelhante a RESTRICT, com verificação em outro momento \\
\hline
\end{tabular}
\end{center}

\section{Escolha orientada por regra de negócio}

A escolha entre \texttt{CASCADE}, \texttt{SET NULL} e \texttt{RESTRICT} não é meramente técnica. Ela depende da semântica da aplicação.

\subsection{Exemplo acadêmico}

Se uma disciplina é removida do catálogo histórico, as matrículas antigas devem desaparecer? Em muitos casos, não. Isso sugere cautela com \texttt{CASCADE}.

\subsection{Exemplo comercial}

Se um cliente for removido, seus pedidos históricos devem desaparecer? Em contextos de auditoria e contabilidade, provavelmente não.

\subsection{Exemplo operacional}

Se uma associação for opcional e transitória, \texttt{SET NULL} pode ser apropriado.

\section{Exemplos práticos ampliados}

\subsection{Alterando estrutura de aluno}

\begin{sqlcode}
ALTER TABLE alunos
ADD COLUMN cpf CHAR(11),
ADD COLUMN status VARCHAR(20) DEFAULT 'ATIVO';
\end{sqlcode}

\subsection{Aplicando unicidade ao CPF}

\begin{sqlcode}
ALTER TABLE alunos
ADD CONSTRAINT uq_cpf UNIQUE (cpf);
\end{sqlcode}

\subsection{Impondo verificação de idade}

\begin{sqlcode}
ALTER TABLE alunos
ADD CONSTRAINT chk_idade_minima
CHECK (idade >= 16);
\end{sqlcode}

\subsection{Criando tabela de professor}

\begin{sqlcode}
CREATE TABLE professor (
    professor_id SERIAL PRIMARY KEY,
    nome VARCHAR(120) NOT NULL,
    titulacao VARCHAR(50) NOT NULL
);
\end{sqlcode}

\subsection{Criando turma com professor opcional}

\begin{sqlcode}
CREATE TABLE turma (
    turma_id SERIAL PRIMARY KEY,
    nome VARCHAR(50) NOT NULL,
    professor_id INTEGER,
    FOREIGN KEY (professor_id)
        REFERENCES professor(professor_id)
        ON DELETE SET NULL
        ON UPDATE CASCADE
);
\end{sqlcode}

\subsection{Criando matrícula com exclusão em cascata}

\begin{sqlcode}
CREATE TABLE matricula (
    matricula_id SERIAL PRIMARY KEY,
    aluno_id INTEGER NOT NULL,
    turma_id INTEGER NOT NULL,
    FOREIGN KEY (aluno_id)
        REFERENCES alunos(aluno_id)
        ON DELETE CASCADE,
    FOREIGN KEY (turma_id)
        REFERENCES turma(turma_id)
        ON DELETE CASCADE
);
\end{sqlcode}

\section{Estudo de caso curto}

Considere o cenário real de um ecossistema acadêmico englobando \texttt{alunos}, \texttt{professor}, \texttt{turma} e \texttt{matricula}. Se, por imposição jurídica, o cadastro de um aluno precisar ser integralmente apagado, todas as suas matrículas deverão sumir de forma encadeada (aplicando-se o \texttt{ON DELETE CASCADE} no vínculo de aluno para matrícula). Por outro lado, caso um professor seja demitido e excluído, a turma não pode desaparecer, mas sim figurar como pendente de lecionador; o banco assume isso configurando o \texttt{ON DELETE SET NULL} no relacionamento de professor para turma. E, finalizando a cascata, se o conselho acadêmico decidir extinguir (remover) a própria turma, logicamente todos os boletins atrelados a ela deverão colapsar junto, reforçando novamente o uso do rigoroso \texttt{ON DELETE CASCADE} na ligação da turma com suas respectivas matrículas.

Esse tipo de raciocínio mostra que a integridade referencial é também uma forma de modelagem do comportamento da aplicação.

\section{Síntese da aula}

Nesta aula, a DDL foi aprofundada para além da criação inicial de estruturas. Foram discutidos o comando \texttt{ALTER TABLE}, as restrições \texttt{DEFAULT}, \texttt{CHECK}, \texttt{UNIQUE} e \texttt{NOT NULL}, além do papel da integridade referencial e das cláusulas \texttt{ON DELETE} e \texttt{ON UPDATE}. O objetivo central foi mostrar que implementar um banco de dados não significa apenas criar tabelas, mas também definir regras estruturais que expressem o comportamento do domínio e garantam consistência ao longo do tempo.